\chapter{Pain after Surgery (Postoperative Pain)}

\section*{Definition}
Postoperative pain is acute pain following a surgical procedure. Effective management of postoperative pain is essential for recovery, mobility, and prevention of complications such as atelectasis, thromboembolism, and chronic pain syndromes.

\section*{Pathophysiology}
Surgical incision causes tissue damage that releases inflammatory mediators (prostaglandins, bradykinin, histamine, substance P), which activate nociceptors. Peripheral and central sensitization can amplify pain signals. Uncontrolled pain leads to sympathetic activation, increased metabolic demand, impaired wound healing, and delayed recovery.

\section*{Causes}
\begin{itemize}
\item Surgical incision and tissue trauma
\item Muscle retraction and dissection
\item Inflammatory response to surgery
\item Nerve injury or irritation
\item Edema and swelling
\item Drain sites and tube placement
\item Patient positioning during surgery
\item Immobilization of joints or limbs
\end{itemize}

\section*{When to See a Doctor}
\begin{itemize}
\item Pain at the surgical site
\item Pain rated on a numeric scale (0-10)
\item Pain that limits movement, coughing, or deep breathing
\item Need for pain medication
\item Pain accompanied by nausea, anxiety, or insomnia
\item Signs of complications: fever, wound drainage, erythema (infection)
\item Pain not responding to prescribed analgesics
\end{itemize}

\section*{Related Symptoms}
\begin{itemize}
\item Acute pain
\item Nausea and vomiting (opioid side effects)
\item Constipation (opioid-induced)
\item Fatigue
\item Anxiety
\item Wound infection
\item Chronic post-surgical pain
\end{itemize}

\section*{Self-Care/Home Management}
\begin{itemize}
\item Take pain medications as prescribed, on schedule
\item Apply ice packs to the surgical site for the first 24-48 hours
\item Elevate the surgical site when possible
\item Use relaxation techniques and deep breathing
\item Move and ambulate as tolerated (promotes recovery)
\item Do not drive or operate machinery while on opioids
\item Report uncontrolled pain to the healthcare provider
\end{itemize}

\section*{How Your Doctor Will Evaluate You}
\begin{itemize}
\item Pain assessment using numeric or visual analog scale
\item Clinical evaluation of the surgical site
\item Vital sign monitoring (pain can cause tachycardia and hypertension)
\end{itemize}

\section*{Treatment Options}
\begin{itemize}
\item Multimodal analgesia: combination of acetaminophen, NSAIDs, and opioids
\item Regional anesthesia: nerve blocks, epidural analgesia
\item Patient-controlled analgesia (PCA) pumps
\item Local anesthetic infiltration at the surgical site
\item Non-pharmacologic: ice, elevation, relaxation, distraction
\item Treat side effects (antiemetics for nausea, stool softeners for constipation)
\end{itemize}
